Orbiter can produce graphical output in a variety of formats:
\begin{enumerate}
\item
TikZ / Latex~\cite{tantau:2022}, see Table~\ref{tab:draw:incidence:structure},
\item
Metapost~\cite{metapost2010},
\item
Bitmap files (.bmp)~\cite{bitmap}, see Table~\ref{tab:bitmap},
\item
Povray, see Section~\ref{sec:povray},
\item
Gnuplot~\cite{gnuplot}, see Section~\ref{sec:gnuplot:interface}.
\end{enumerate}


\bigskip



The poset classification algorithm from Sections~\ref{sec:subsetorbits} 
and~\ref{sec:subspaceorbits} computes partially ordered sets. 
The posets are created using the \verb'-draw_poset' option in 
the poset classification control command package, see Table~\ref{tab:pc:options}. 
The posets are stored in a file with extension \verb'.layered_graph'.
These files can be drawn using the \verb'-draw_layered_graph' command.
The commands in Table~\ref{tab:draw:options} and Table~\ref{tab:draw:options:b} 
show ways in which to customize the drawings.
\orbitertableofcommands{layered_graph_draw_options_1}
\orbitertableofcommands{layered_graph_draw_options_2}
Let us consider an example. 
Suppose we are interested in the Schreier trees 
of a permutation group represented in a stabilizer chain.
We take $\PGL(4,2)$ in its action on the wedge product.
The command

\sourcecode{PGL_6_2_Wedge_4_2_detached_report.tex}

produces a report.
The following tree is taken from this report. It represents the  
first basic orbit 
in the stabilizer chain of the group in that action.

\orbiteroutput{GRAPHICS/LATEX/orbits_PGL42_wedge_detached}


\bigskip



The command 

\sourcecode{schreier_tree_graphical_output}

draws the $6$th Schreier tree in the classification of orbits of $\PGL(4,2)$ on 
homogeneous polynomials of degree 3 in 4 variables, 
shown below

\orbiteroutput{GRAPHICS/WRAPPERS/on_poly_schreier_tree_6_draw}

The orbit has length $420.$ 

\bigskip


